\chapter{Detailed implementation plan}
\label{app:implementation}
\label{app:code}
\chapterpromise{A module-by-module plan for implementing the stated probability models, numerical methods, experiments, data schemas, tests, and reproducibility records.}

\section{Scope and non-goals}
The implementation work preserves the author-approved method: generalized hierarchical Bayesian
segmentation with exponential-family segment models, irregular-design priors, multiple-sequence
hierarchies, and known or latent groups. The core software separates observation-model calculations
from sum-product and max-sum recursions over ordered partitions. The plan does not introduce an
online detector, reinterpret the finite latent-group criterion as a normalized mixture likelihood,
change the scientific endpoint, or overwrite archived numerical results. A corrected computation
receives a new result identifier and records its relationship to the earlier output.

\section{Repository organization}
\begin{description}
\item[\path{src/bayesbreak/}] Observation families, segment priors, sum-product recursion, max-sum MAP calculation, multiple-sequence models, grouped models, and prediction.
\item[\path{configs/}] Versioned model, comparator, dataset, and execution specifications.
\item[\path{experiments/}] Drivers for \texttt{EPR-BB-001} through \texttt{EPR-BB-015}.
\item[\path{scripts/}] Data-schema, result-metadata, bibliography, build, and release validators.
\item[\path{tests/}] Unit, property-based, numerical, statistical, regression, and reproducibility tests.
\item[\path{data/}] Immutable source records and versioned processed datasets.
\item[\path{artifacts/}] Result files, figures, tables, numerical error records, and records of excluded computations.
\item[\path{docs/schemas/}] Schemas for results, numerical-error summaries, metrics, and reproducibility metadata.
\end{description}
Responsibility is organized by component: observation models, partition inference, posterior
prediction and metrics, experiments and result metadata, and release verification. A task is
complete only when the corresponding mathematical invariant and regression test pass together.

\section{Public functions and returned objects}
\begin{description}
\item[Observation-family calculation] Returns a log segment marginal likelihood, requested
posterior moments or posterior state, support information, and, for a numerical approximation, the
available convergence and error information for each admissible segment.
\item[Partition prior] Returns log segment cohesion and log interior-boundary hazard separately,
including zero-probability cases.
\item[Partition inference] Uses the same prior-adjusted segment scores in sum-product and max-sum
recursions and returns marginal likelihoods, $P(k\mid y)$, changepoint marginals, Bayes curves, and a
joint MAP partition.
\item[Posterior prediction] Dispatches according to the fitted observation family and applies the
stated coordinate-support or extrapolation rule.
\item[Changepoint metrics] Validates coordinate axes and uses maximum-cardinality,
minimum-distance one-to-one matching; matched-location MAE is \NA{} when no pair is matched.
\item[Result record] Distinguishes execution status from validity for a specified analysis and does
not mutate an archived result in place.
\end{description}

\section{Task registry and acceptance criteria}
% Generated from shared/handoffs/coding_agent_handoff.json.
\small
\begin{longtable}{@{}p{0.16\linewidth}p{0.37\linewidth}p{0.37\linewidth}@{}}
\caption{Implementation task registry. Detailed files, interfaces, and invariants are in Appendix~\ref{app:coding-handoff}.}\label{tab:code-tasks}\\
\toprule
\textbf{ID} & \textbf{Work item} & \textbf{Acceptance criterion}\\
\midrule
\endfirsthead
\toprule
\textbf{ID} & \textbf{Work item} & \textbf{Acceptance criterion}\\
\midrule
\endhead
\texttt{CODE-BB-001} & Implement explicit segment-cohesion and boundary-hazard prior specification & Analytic small-n enumeration and sensitivity fixtures pass. \\
\texttt{CODE-BB-002} & Integrate repaired prior with exact DP and property tests & EPR-BB-001 randomized invariant suite passes. \\
\texttt{CODE-BB-003} & Repair the finite latent-group criterion and restart selection & All exit paths, ties, collapsed groups, and restart ordering tests pass. \\
\texttt{CODE-BB-004} & Harden grouped and replicate log-domain aggregation & Adversarial magnitude tests pass without overflow/underflow-induced ranking changes. \\
\texttt{CODE-BB-005} & Add a schema for nonconjugate segment-error assessment & Nested tolerances tighten the verified bound or return an explicit failure state; synthetic propagation checks hold. \\
\texttt{CODE-BB-006} & Implement Beta-observation posterior predictive scoring & Normalization/reference tests and edge-case finite checks pass. \\
\texttt{CODE-BB-007} & Make extrapolation policy explicit & Support-boundary and provenance tests pass for every policy. \\
\texttt{CODE-BB-008} & Canonicalize boundary matching and metrics & Hand-computed, tie, empty, and symmetry fixtures pass. \\
\texttt{CODE-BB-009} & Validate comparator axes and tuning budgets & CGH cached single-trace artifact is rejected deterministically; valid raw matrix route is covered. \\
\texttt{CODE-BB-010} & Version result and provenance schemas & Every new sidecar validates and every archived sidecar has a documented migration/read-only path. \\
\texttt{CODE-BB-011} & Synchronize real-data loaders, cards, captions, and hashes & Well-log, CGH, SPX, and methylKit cards validate against generated sidecars and captions. \\
\texttt{CODE-BB-012} & Profile and re-run the full test suite & All required tests terminate under pinned profiles or remain explicitly blocked with owner and reason. \\
\texttt{CODE-BB-013} & Reconcile version, CI, changelog, and preprint lineage & Clean checkout reports one version and citation target; lineage table is complete. \\
\texttt{CODE-BB-014} & Render canonical handoff and appendix fragments & Book appendix and task registry pass semantic and ID synchronization checks. \\
\texttt{CODE-BB-015} & Regenerate figures and tables with semantic validation & Every figure/table has source hashes, status metadata, and caption anchor; visual QA passes. \\
\texttt{CODE-BB-016} & Automate bibliography annotation completeness & Completeness checker returns zero missing, duplicate, orphaned, or relationship-free entries. \\
\bottomrule
\end{longtable}
\normalsize


\section{Data, result metadata, and leakage controls}
Each dataset receives a source record, immutable hash, license or access note, raw schema, and a
versioned sequence of transformations. Splits are fixed before fitting. Simulation truth is stored
separately from fitted outputs. Comparator tuning uses only the declared training or validation
information. Agreement with a \BayesBreak fit is not relabelled as external truth. The well-log,
array-CGH, equity, and methylation datasets retain their executed preprocessing and hashes; a
corrected or extended analysis is stored as a new version.

\section{Execution, caching, and failure states}
A run records code, data, model specification, environment, random seed, wall time, memory,
convergence state, and output hashes. Cached segment calculations are reused only when the
observation family, descriptors, priors, support, and code hash agree. A nonconjugate method may
finish with a numerical value but without a verified error bound; a comparator calculation may
finish but be invalid for a specified comparison; a test may remain unresolved because it does not
terminate under the execution limit. These states are reported separately.

\section{Tests}
Unit tests cover sufficient statistics, conjugate normalizers, support, and segment factorization.
Property-based tests compare sum-product and max-sum recursions with exhaustive enumeration on
small problems. Numerical tests check log-domain normalization and segmentwise error calculations.
Statistical tests cover calibration, null simulations, and repeated-seed behavior. Regression tests
preserve archived result identifiers and ensure that excluded values cannot enter incompatible
figures or tables. Reproducibility tests rebuild figures and tables and compare declared hashes or
numerical tolerances.

\section{Implementation order and stopping criteria}
The dependency order is: segment-cohesion and boundary-hazard semantics; dynamic-programming
invariants; latent-group and shared-changepoint numerical corrections; posterior prediction and
changepoint metrics; comparator and data schemas; full-suite profiling; controlled recomputation;
figure and table regeneration; and source synchronization. Each stage has a pre-specified
acceptance criterion. When a theorem assumption, numerical error bound, comparator definition, or
test cannot be satisfied, the corresponding claim is narrowed or the calculation remains excluded
from that scientific interpretation. An unimplemented item remains explicitly incomplete.
